% =============================================================================
% Chapter 4 — BASIC Program Space ($0280–$9FFF)
% =============================================================================
\chapter{BASIC Program Space}

\epigraph{``Forty kilobytes ought to be enough for anybody's BASIC
program.''}{\textit{--- Apocryphal, ca.\ 1982}}

\section*{Introduction}

The roughly 40\,KB region from \$0280 to \$9FFF is where BASIC does all of its
living.  Every line you type, every variable you assign, every array you
dimension, and every string you concatenate ends up somewhere in this space.
Understanding how the interpreter organises its data here is essential for
writing programs that are memory-efficient, for debugging mysterious
\texttt{?OUT OF MEMORY} errors, and for writing machine-language routines that
can inspect or manipulate BASIC's internal state.

The region is not one flat pool.  It is divided into four distinct areas ---
program text, simple variables, array storage, and the string heap --- whose
boundaries shift dynamically as your program runs.  Six zero-page pointer
pairs (documented in Chapter~1) track where each area begins and ends.

% ═══════════════════════════════════════════════════════════════════════════════
\section{Memory Layout Overview}
% ═══════════════════════════════════════════════════════════════════════════════

\memrange{0280}{9FFF}{BASIC Program Space}{R/W}{0}

The entire region is ordinary read/write RAM.  At boot, EhBASIC sets the
program-start pointer (\addr{0079}/\addr{007A}) to \$0280 and the
end-of-memory pointer (\addr{0085}/\addr{0086}) to \$A000.  Everything in
between is available for your program and its data.

The layout from low addresses to high is:

\begin{center}
\begin{tikzpicture}[
  block/.style={draw, minimum width=18em, minimum height=2.2em,
                font=\ttfamily\small, anchor=north},
  lbl/.style={font=\ttfamily\scriptsize, text=retrogray, anchor=east},
  ptr/.style={font=\ttfamily\scriptsize\bfseries, text=retroblue, anchor=west}
]
  % Blocks (top to bottom)
  \node[block, fill=retropale] (prog) at (0,0)
    {Tokenised Program};
  \node[block, fill=retrocyan!12!retrocream, below=0pt of prog] (vars)
    {Simple Variables};
  \node[block, fill=retrogreen!12!retrocream, below=0pt of vars] (arrs)
    {Array Storage};
  \node[block, fill=white, below=0pt of arrs, minimum height=4em] (free)
    {Free Memory};
  \node[block, fill=retroorange!15!retrocream, below=0pt of free] (str)
    {String Heap (grows \(\downarrow\))};

  % Left-side address labels
  \node[lbl] at ([xshift=-1em]prog.north west) {\$0280};
  \node[lbl] at ([xshift=-1em]str.south west)  {\$9FFF};

  % Right-side pointer labels
  \node[ptr] at ([xshift=1em]prog.north east)
    {\(\leftarrow\) Smeml/h (\$79/\$7A)};
  \node[ptr] at ([xshift=1em]vars.north east)
    {\(\leftarrow\) Svarl/h (\$7B/\$7C)};
  \node[ptr] at ([xshift=1em]arrs.north east)
    {\(\leftarrow\) Sarryl/h (\$7D/\$7E)};
  \node[ptr] at ([xshift=1em]free.north east)
    {\(\leftarrow\) Earryl/h (\$7F/\$80)};
  \node[ptr] at ([xshift=1em]str.north east)
    {\(\leftarrow\) Sstorl/h (\$81/\$82)};
  \node[ptr] at ([xshift=1em]str.south east)
    {\(\leftarrow\) Ememl/h (\$85/\$86)};
\end{tikzpicture}
\end{center}

The invariant is always:
\begin{center}
\texttt{Smem} $\leq$ \texttt{Svar} $\leq$ \texttt{Sarry} $\leq$
\texttt{Earry} $\leq$ \texttt{Sstor} $\leq$ \texttt{Emem}
\end{center}

Program text, variables, and arrays grow \emph{upward} from \$0280.  The
string heap grows \emph{downward} from \$9FFF.  When \texttt{Earry} meets
\texttt{Sstor}, memory is exhausted and EhBASIC raises
\texttt{?OUT OF MEMORY}.

\begin{tipbox}
You can watch these pointers in real time:
\begin{lstlisting}
10 PRINT "PROG  :$";HEX$(DEEK($79))
20 PRINT "VARS  :$";HEX$(DEEK($7B))
30 PRINT "ARRAYS:$";HEX$(DEEK($7D))
40 PRINT "EARRY :$";HEX$(DEEK($7F))
50 PRINT "SSTOR :$";HEX$(DEEK($81))
60 PRINT "EMEM  :$";HEX$(DEEK($85))
\end{lstlisting}
\end{tipbox}

% ═══════════════════════════════════════════════════════════════════════════════
\section{Tokenised Line Format}
% ═══════════════════════════════════════════════════════════════════════════════

BASIC programs are not stored as ASCII text.  When you enter a line, EhBASIC
\emph{tokenises} it --- replacing keywords with single-byte (or two-byte)
tokens, while leaving literal text, numbers, and punctuation untouched.  This
saves memory and speeds up the interpreter, which can dispatch on a single
byte rather than parsing multi-character keywords.

Each line in memory has the following structure:

\begin{center}
\begin{tabular}{clp{24em}}
\toprule
\textbf{Offset} & \textbf{Size} & \textbf{Description} \\
\midrule
+0 & 2 bytes & Pointer to next line (little-endian).
               \$0000 marks end of program. \\
+2 & 2 bytes & Line number (little-endian, 0--65535). \\
+4 & $n$ bytes & Tokenised content. \\
+4+$n$ & 1 byte & \$00 terminator. \\
\bottomrule
\end{tabular}
\end{center}

Lines form a singly linked list.  The interpreter walks the list by reading
the next-line pointer; when it finds \$0000, the program is over.  Line
numbers need not be contiguous --- \texttt{10}, \texttt{20}, \texttt{100} is
perfectly valid.

\subsection{Worked Example: \texttt{10 PRINT "HELLO"}}

Suppose this one-line program is stored beginning at \$0280.  Here is the
byte-by-byte hex dump:

\begin{center}
\begin{tabular}{ccl}
\toprule
\textbf{Address} & \textbf{Hex} & \textbf{Meaning} \\
\midrule
\$0280 & \texttt{8D 02} & Next-line pointer = \$028D \\
\$0282 & \texttt{0A 00} & Line number = \$000A (10) \\
\$0284 & \texttt{9F}    & PRINT token \\
\$0285 & \texttt{20}    & Space (literal) \\
\$0286 & \texttt{22}    & \texttt{"} (open quote) \\
\$0287 & \texttt{48}    & \texttt{H} \\
\$0288 & \texttt{45}    & \texttt{E} \\
\$0289 & \texttt{4C}    & \texttt{L} \\
\$028A & \texttt{4C}    & \texttt{L} \\
\$028B & \texttt{4F}    & \texttt{O} \\
\$028C & \texttt{00}    & End-of-line terminator \\
\midrule
\$028D & \texttt{00 00} & Next-line pointer = \$0000 (end of program) \\
\$028F & \texttt{00 00} & (padding / end marker) \\
\bottomrule
\end{tabular}
\end{center}

\begin{notebox}
The closing quote character (\texttt{\$22}) is not stored --- EhBASIC strips
trailing quotes at the end of \texttt{PRINT} arguments.  The string literal
ends at the \$00 line terminator.
\end{notebox}

% ═══════════════════════════════════════════════════════════════════════════════
\section{Token Table}
\label{sec:token-table}
% ═══════════════════════════════════════════════════════════════════════════════

NovaBASIC uses two token spaces.  \textbf{Primary tokens} occupy the range
\$80--\$FF (one byte each).  \textbf{Extended tokens} are two bytes: a prefix
byte \$01 followed by a sub-token byte.

\subsection{Primary Tokens (\$80--\$FF)}

\subsubsection*{Command Tokens}

\begin{center}
\begin{longtable}{cl|cl|cl}
\toprule
\textbf{Hex} & \textbf{Keyword} &
\textbf{Hex} & \textbf{Keyword} &
\textbf{Hex} & \textbf{Keyword} \\
\midrule
\endhead
\$80 & END      & \$90 & RETURN   & \$A0 & CONT     \\
\$81 & FOR      & \$91 & REM      & \$A1 & LIST     \\
\$82 & NEXT     & \$92 & STOP     & \$A2 & CLEAR    \\
\$83 & DATA     & \$93 & ON       & \$A3 & NEW      \\
\$84 & INPUT    & \$94 & NULL     & \$A4 & WIDTH    \\
\$85 & DIM      & \$95 & INC      & \$A5 & GET      \\
\$86 & READ     & \$96 & WAIT     & \$A6 & SWAP     \\
\$87 & LET      & \$97 & LOAD     & \$A7 & BITSET   \\
\$88 & DEC      & \$98 & SAVE     & \$A8 & BITCLR   \\
\$89 & GOTO     & \$99 & DEF      & \$A9 & IRQ      \\
\$8A & RUN      & \$9A & POKE     & \$AA & NMI      \\
\$8B & IF       & \$9B & DOKE     & \$AB & CLS      \\
\$8C & RESTORE  & \$9C & SYS      & \$AC & COLOR    \\
\$8D & GOSUB    & \$9D & DO       & \$AD & LOCATE   \\
\$8E & RETIRQ   & \$9E & LOOP     & \$AE & PLOT     \\
\$8F & RETNMI   & \$9F & PRINT    & \$AF & UNPLOT   \\
\bottomrule
\end{longtable}
\end{center}

\begin{center}
\begin{longtable}{cl|cl|cl}
\toprule
\textbf{Hex} & \textbf{Keyword} &
\textbf{Hex} & \textbf{Keyword} &
\textbf{Hex} & \textbf{Keyword} \\
\midrule
\endhead
\$B0 & LINE        & \$B8 & SPRITE      & \$C0 & VSYNC    \\
\$B1 & CIRCLE      & \$B9 & SPRITESHAPE & \$C1 & TAB(     \\
\$B2 & RECT        & \$BA & SPRITESET   & \$C2 & ELSE     \\
\$B3 & FILL        & \$BB & SPRITEDATA  & \$C3 & TO       \\
\$B4 & PAINT       & \$BC & SOUND       & \$C4 & FN       \\
\$B5 & MODE        & \$BD & VOLUME      & \$C5 & SPC(     \\
\$B6 & GCLS        & \$BE & INSTRUMENT  & \$C6 & THEN     \\
\$B7 & GCOLOR      & \$BF & WAVE        & \$C7 & NOT      \\
\bottomrule
\end{longtable}
\end{center}

\subsubsection*{Operator and Structure Tokens}

\begin{center}
\begin{tabular}{cl|cl|cl}
\toprule
\textbf{Hex} & \textbf{Token} &
\textbf{Hex} & \textbf{Token} &
\textbf{Hex} & \textbf{Token} \\
\midrule
\$C8 & STEP     & \$CC & +     & \$D0 & \^{}     \\
\$C9 & UNTIL    & \$CD & -     & \$D1 & AND      \\
\$CA & WHILE    & \$CE & *     & \$D2 & EOR      \\
\$CB & OFF      & \$CF & /     & \$D3 & OR       \\
     &          & \$D4 & >>    & \$D6 & >        \\
     &          & \$D5 & <<    & \$D7 & =        \\
     &          &      &       & \$D8 & <        \\
\bottomrule
\end{tabular}
\end{center}

\subsubsection*{Function Tokens}

\begin{center}
\begin{longtable}{cl|cl|cl}
\toprule
\textbf{Hex} & \textbf{Function} &
\textbf{Hex} & \textbf{Function} &
\textbf{Hex} & \textbf{Function} \\
\midrule
\endhead
\$D9 & SGN(       & \$E3 & BITTST(    & \$ED & SPRITEX(     \\
\$DA & INT(       & \$E4 & MAX(       & \$EE & SPRITEY(     \\
\$DB & ABS(       & \$E5 & MIN(       & \$EF & COLLISION(   \\
\$DC & USR(       & \$E6 & PI         & \$F0 & BUMPED(      \\
\$DD & FRE(       & \$E7 & TWOPI      &      &              \\
\$DE & POS(       & \$E8 & VARPTR(    &      &              \\
\$DF & SQR(       & \$E9 & LEFT\$(    &      &              \\
\$E0 & RND(       & \$EA & RIGHT\$(   &      &              \\
\$E1 & LOG(       & \$EB & MID\$(     &      &              \\
\$E2 & EXP(       & \$EC & --- (reserved) & &              \\
\midrule
\multicolumn{6}{l}{\small Continued:} \\
\midrule
\$E3 & COS(       & \$E7 & PEEK(      & \$EB & STR\$(     \\
\$E4 & SIN(       & \$E8 & DEEK(      & \$EC & VAL(       \\
\$E5 & TAN(       & \$E9 & SADD(      & \$ED & ASC(       \\
\$E6 & ATN(       & \$EA & LEN(       & \$EE & UCASE\$(   \\
\bottomrule
\end{longtable}
\end{center}

\begin{warningbox}
The function token table above is presented in the order matching the
\texttt{tokens.json} source file.  Here is the complete mapping from
\$D9 to \$FF:
\end{warningbox}

\begin{center}
\begin{longtable}{cl|cl}
\toprule
\textbf{Hex} & \textbf{Function} &
\textbf{Hex} & \textbf{Function} \\
\midrule
\endhead
\$D9 & SGN(       & \$E9 & SADD(      \\
\$DA & INT(       & \$EA & LEN(       \\
\$DB & ABS(       & \$EB & STR\$(     \\
\$DC & USR(       & \$EC & VAL(       \\
\$DD & FRE(       & \$ED & ASC(       \\
\$DE & POS(       & \$EE & UCASE\$(   \\
\$DF & SQR(       & \$EF & LCASE\$(   \\
\$E0 & RND(       & \$F0 & CHR\$(     \\
\$E1 & LOG(       & \$F1 & HEX\$(     \\
\$E2 & EXP(       & \$F2 & BIN\$(     \\
\$E3 & COS(       & \$F3 & BITTST(    \\
\$E4 & SIN(       & \$F4 & MAX(       \\
\$E5 & TAN(       & \$F5 & MIN(       \\
\$E6 & ATN(       & \$F6 & PI         \\
\$E7 & PEEK(      & \$F7 & TWOPI      \\
\$E8 & DEEK(      & \$F8 & VARPTR(    \\
     &            & \$F9 & LEFT\$(    \\
     &            & \$FA & RIGHT\$(   \\
     &            & \$FB & MID\$(     \\
     &            & \$FC & SPRITEX(   \\
     &            & \$FD & SPRITEY(   \\
     &            & \$FE & COLLISION( \\
     &            & \$FF & BUMPED(    \\
\bottomrule
\end{longtable}
\end{center}

\subsection{Extended Tokens (\$01 + sub-token)}

When the tokeniser encounters a keyword that does not fit in the primary
table, it emits the prefix byte \$01 followed by a sub-token.  These are the
NovaVM extensions --- commands and functions not present in the original
EhBASIC.

\begin{center}
\begin{longtable}{ccl|ccl}
\toprule
\textbf{Prefix} & \textbf{Sub} & \textbf{Keyword} &
\textbf{Prefix} & \textbf{Sub} & \textbf{Keyword} \\
\midrule
\endhead
\$01 & \$01 & DIR         & \$01 & \$19 & NOPEN       \\
\$01 & \$02 & DEL         & \$01 & \$1A & NCLOSE      \\
\$01 & \$03 & XMEM        & \$01 & \$1B & NLISTEN     \\
\$01 & \$04 & XBANK       & \$01 & \$1C & NACCEPT     \\
\$01 & \$05 & XPOKE       & \$01 & \$1D & NSEND       \\
\$01 & \$06 & XPEEK(      & \$01 & \$1E & NRECV\$(    \\
\$01 & \$07 & STASH       & \$01 & \$1F & NSTATUS(    \\
\$01 & \$08 & FETCH       & \$01 & \$21 & NREADY(     \\
\$01 & \$09 & XFREE       & \$01 & \$22 & NLEN        \\
\$01 & \$0A & XRESET      & \$01 & \$23 & DMACOPY     \\
\$01 & \$0B & XALLOC      & \$01 & \$24 & DMAFILL     \\
\$01 & \$0C & XDIR        & \$01 & \$25 & DMASTATUS   \\
\$01 & \$0D & XDEL        & \$01 & \$27 & DMAERR      \\
\$01 & \$0E & XMAP        & \$01 & \$28 & DMACOUNT    \\
\$01 & \$0F & XUNMAP      & \$01 & \$29 & MIDPLAY     \\
\$01 & \$10 & GSAVE       & \$01 & \$2A & MIDSTOP     \\
\$01 & \$11 & GLOAD       & \$01 & \$2B & DIROPEN     \\
\$01 & \$12 & SIDPLAY     & \$01 & \$2C & DIRNEXT     \\
\$01 & \$13 & SIDSTOP     & \$01 & \$2D & DIRNAM\$(   \\
\$01 & \$14 & MUSIC       & \$01 & \$2E & DIRSIZ      \\
\$01 & \$15 & PLAYING     & \$01 & \$2F & DIRTYP      \\
\$01 & \$16 & MNOTE(      & \$01 & \$30 & META\$(     \\
\$01 & \$17 & COPPER      & & & \\
\$01 & \$18 & RESET       & & & \\
\midrule
\$01 & \$40 & BLITCOPY    & \$01 & \$46 & HELP        \\
\$01 & \$41 & BLITFILL    & \$01 & \$47 & NCC         \\
\$01 & \$42 & BLITSTATUS  & \$01 & \$48 & FONT        \\
\$01 & \$43 & BLITERR     & \$01 & \$49 & CD          \\
\$01 & \$44 & BLITCOUNT   & \$01 & \$4A & MKDIR       \\
\$01 & \$45 & BITTGL      & \$01 & \$4B & RMDIR       \\
     &      &             & \$01 & \$4C & FORMAT      \\
     &      &             & \$01 & \$4D & MOUNT       \\
     &      &             & \$01 & \$4E & UNMOUNT     \\
     &      &             & \$01 & \$4F & PWD         \\
\bottomrule
\end{longtable}
\end{center}

\begin{notebox}
Inside quoted strings, no tokenisation occurs.  The bytes between double
quotes are stored as literal ASCII.  This is why you can safely embed keyword
substrings in strings: \texttt{"FORGET"} does not become \texttt{"}\,\$81\,%
\texttt{GET"}.
\end{notebox}

% ═══════════════════════════════════════════════════════════════════════════════
\section{Simple Variable Storage}
% ═══════════════════════════════════════════════════════════════════════════════

Immediately after the last program line (at the address in
\addr{007B}/\addr{007C}), EhBASIC stores simple (non-array) variables.  Each
variable occupies either six or seven bytes, depending on its type.

\subsection{Variable Name Encoding}

The first two bytes of each variable entry encode the name:

\begin{center}
\begin{tabular}{cp{28em}}
\toprule
\textbf{Byte} & \textbf{Contents} \\
\midrule
Byte 1 & First significant character of the name (A--Z).
         Bit~7 set indicates a string variable (\texttt{\$} suffix). \\
Byte 2 & Second significant character, or \$00 if the name is a single
         letter.  Bit~7 set indicates an integer function variable. \\
\bottomrule
\end{tabular}
\end{center}

EhBASIC only distinguishes variables by their first two characters.
\texttt{SCORE} and \texttt{SCREEN} are the same variable.

\subsection{Numeric Variables}

A numeric variable entry is six bytes:

\begin{center}
\begin{tabular}{cl}
\toprule
\textbf{Offset} & \textbf{Contents} \\
\midrule
+0 & Name byte 1 (bit 7 clear = numeric) \\
+1 & Name byte 2 \\
+2..+5 & 4-byte floating-point value (EhBASIC internal format) \\
\bottomrule
\end{tabular}
\end{center}

\begin{notebox}
EhBASIC uses a non-standard 32-bit floating-point format: 1 byte exponent
(biased by 128), followed by a 24-bit mantissa with an implied leading 1 bit.
The sign is stored in the high bit of the first mantissa byte.  This gives
roughly 6.5 decimal digits of precision.
\end{notebox}

\subsection{String Variables}

A string variable entry is also six bytes, but the value portion is
different:

\begin{center}
\begin{tabular}{cl}
\toprule
\textbf{Offset} & \textbf{Contents} \\
\midrule
+0 & Name byte 1 (bit 7 set = string) \\
+1 & Name byte 2 \\
+2 & String length (0--255) \\
+3 & (unused / padding) \\
+4..+5 & Pointer to string data in the string heap (little-endian) \\
\bottomrule
\end{tabular}
\end{center}

The string data itself lives in the string heap at the top of memory.  The
variable entry only holds a descriptor pointing to it.

% ═══════════════════════════════════════════════════════════════════════════════
\section{Array Storage}
% ═══════════════════════════════════════════════════════════════════════════════

Arrays are stored after simple variables, between \addr{007D}/\addr{007E} and
\addr{007F}/\addr{0080}.  Each array entry consists of a header followed by
packed element values.

\subsection{Array Header}

\begin{center}
\begin{tabular}{cl}
\toprule
\textbf{Offset} & \textbf{Contents} \\
\midrule
+0 & Name byte 1 (same encoding as simple variables) \\
+1 & Name byte 2 \\
+2..+3 & Total size of this array entry in bytes (little-endian),
         including the header \\
+4 & Number of dimensions (1--255) \\
+5..+$n$ & 2 bytes per dimension: element count (little-endian).
           \texttt{DIM A(9)} stores 10 (0 through 9). \\
\bottomrule
\end{tabular}
\end{center}

\subsection{Array Element Data}

After the dimension sizes, elements are stored in row-major order:

\begin{itemize}
  \item \textbf{Numeric arrays}: 4 bytes per element (same floating-point
        format as simple variables).
  \item \textbf{String arrays}: 3 bytes per element (1 byte length + 2 bytes
        pointer into the string heap).
\end{itemize}

\begin{warningbox}
Arrays consume memory rapidly.  \texttt{DIM A(100)} allocates a header plus
$101 \times 4 = 404$ bytes for numeric data.  A two-dimensional array
\texttt{DIM B(10,10)} needs $11 \times 11 \times 4 = 484$ bytes.  Check
\texttt{FRE(0)} before dimensioning large arrays.
\end{warningbox}

% ═══════════════════════════════════════════════════════════════════════════════
\section{String Heap}
% ═══════════════════════════════════════════════════════════════════════════════

String data occupies the top of the BASIC program space, growing
\emph{downward} from \addr{0085}/\addr{0086} (Ememl/h).  The current bottom
of the heap is tracked by \addr{0081}/\addr{0082} (Sstorl/h).

When you assign \texttt{A\$="HELLO"}, EhBASIC:
\begin{enumerate}
  \item Decrements Sstor by the string length (5).
  \item Copies the five bytes into the new location.
  \item Updates the variable entry's pointer to the new address.
\end{enumerate}

\subsection{Garbage Collection}

Over time, reassigning strings leaves orphaned data in the heap.  When free
memory runs low, EhBASIC performs \emph{garbage collection}: it scans all
string variables and array elements, compacts the live strings toward the top
of memory, and reclaims the dead space.

This can cause a noticeable pause --- the dreaded ``BASIC hiccup'' familiar
to 8-bit programmers.  Programs that do heavy string manipulation may
experience periodic delays.

\begin{tipbox}
You can force garbage collection at a convenient time by calling
\texttt{FRE("")}.  This is useful before a timing-critical section:
\begin{lstlisting}
1000 X=FRE("") : REM compact strings now
1010 REM -- timing-sensitive code follows --
\end{lstlisting}
\end{tipbox}

% ═══════════════════════════════════════════════════════════════════════════════
\section{Free Memory}
% ═══════════════════════════════════════════════════════════════════════════════

The free memory region is the gap between the end of array storage
(\addr{007F}/\addr{0080}) and the bottom of the string heap
(\addr{0081}/\addr{0082}).  The \texttt{FRE(0)} function returns the size of
this gap in bytes.

\begin{lstlisting}
10 PRINT "FREE MEMORY:"; FRE(0); "BYTES"
20 DIM A(100)
30 PRINT "AFTER DIM:"; FRE(0); "BYTES"
\end{lstlisting}

On a freshly booted NovaVM with no program loaded, free memory is
approximately 40,192 bytes (\$9D00 in hex) --- the full span from \$0280 to
\$9FFF minus the few bytes EhBASIC uses for the end-of-program marker.

\begin{notebox}
\texttt{FRE(0)} returns a numeric result.  \texttt{FRE("")} also triggers
garbage collection and then returns the free count.  Always use
\texttt{FRE("")} if you want an accurate count after heavy string use.
\end{notebox}

% ═══════════════════════════════════════════════════════════════════════════════
\section{Walking the Data Structures}
% ═══════════════════════════════════════════════════════════════════════════════

The linked-list nature of tokenised program lines and the packed variable
tables can be explored directly from BASIC using \texttt{PEEK} and
\texttt{DEEK}.

\subsection{Walking the Program Line List}

This program follows the next-line pointers from the start of the program,
printing each line number it finds:

\begin{lstlisting}
1000 REM -- Walk program lines --
1010 P=DEEK($79) : REM start of program
1020 NP=DEEK(P) : REM next-line pointer
1030 IF NP=0 THEN PRINT "END OF PROGRAM" : END
1040 LN=DEEK(P+2) : REM line number
1050 SZ=NP-P : REM line size in bytes
1060 PRINT "LINE";LN;"AT $";HEX$(P);"SIZE";SZ
1070 P=NP : GOTO 1020
\end{lstlisting}

\subsection{Dumping a Tokenised Line}

This program dumps the raw bytes of a specific line, showing how tokens
replace keywords:

\begin{lstlisting}
2000 REM -- Hex dump of a BASIC line --
2010 INPUT "LINE NUMBER";TL
2020 P=DEEK($79)
2030 NP=DEEK(P) : IF NP=0 THEN PRINT "NOT FOUND" : END
2040 LN=DEEK(P+2) : IF LN=TL THEN 2070
2050 P=NP : GOTO 2030
2060 REM -- Found it --
2070 SZ=NP-P
2080 PRINT "LINE";TL;"AT $";HEX$(P);"-";
2090 PRINT "$";HEX$(P+SZ-1);" (";SZ;"BYTES)"
2100 FOR I=0 TO SZ-1
2110 PRINT HEX$(PEEK(P+I));" ";
2120 IF (I AND 15)=15 THEN PRINT
2130 NEXT I
2140 PRINT
\end{lstlisting}

\subsection{Showing the Memory Map}

A quick-reference program that displays all six pointer pairs and the
computed region sizes:

\begin{lstlisting}
3000 REM -- BASIC memory map --
3010 SM=DEEK($79):SV=DEEK($7B)
3020 SA=DEEK($7D):EA=DEEK($7F)
3030 SS=DEEK($81):EM=DEEK($85)
3040 PRINT "PROGRAM  : $";HEX$(SM);" - $";HEX$(SV-1)
3050 PRINT "          ";SV-SM;"BYTES"
3060 PRINT "VARIABLES: $";HEX$(SV);" - $";HEX$(SA-1)
3070 PRINT "          ";SA-SV;"BYTES"
3080 PRINT "ARRAYS   : $";HEX$(SA);" - $";HEX$(EA-1)
3090 PRINT "          ";EA-SA;"BYTES"
3100 PRINT "FREE     : $";HEX$(EA);" - $";HEX$(SS-1)
3110 PRINT "          ";SS-EA;"BYTES"
3120 PRINT "STRINGS  : $";HEX$(SS);" - $";HEX$(EM-1)
3130 PRINT "          ";EM-SS;"BYTES"
3140 PRINT
3150 PRINT "FRE(0) = ";FRE(0)
\end{lstlisting}
